% !TeX root = ../main.tex
\section{Introduction}

% 中文：由于源域和目标域之间的分布差异，预训练模型在部署阶段往往会出现性能下降。
% 持续测试时适配（CTTA）通过在推理过程中利用在线无标签测试数据更新模型参数来缓解这一问题，
% 因而对于真实世界的视觉理解和嵌入式智能等应用具有重要意义。
% 近期研究进一步将 CTTA 扩展到多模态场景（MM-CTTA），例如音频-视频动作识别和事件分类，
% 并通过模态可靠性建模、模态选择、非对称适配或监督修正取得了较强性能。
% 然而，现有 MM-CTTA 方法大多仍依赖反向传播，并通常通过熵最小化等辅助目标更新模型参数。
% 这与真实端侧部署存在明显矛盾：端侧设备通常只有有限的计算资源和内存用于在线适配，
% 或者部署在不支持反向传播的推理引擎和量化模型组件上。
% 因此，本文研究 MM-CTTA 在端侧场景下的部署问题，即在资源受限、仅支持前向传播、单样本流式输入的条件下实现稳定有效的多模态持续测试时适配。
Due to source-target distribution changes, pretrained models often suffer severe performance degradation after deployment.
Continual test-time adaptation (CTTA) alleviates this issue by updating model parameters during inference with online unlabeled test data, making it important for real-world applications such as visual understanding and embedded intelligence \citep{wang2022continualtta}.
Recent studies have further extended CTTA to multimodal settings (MM-CTTA), such as audio-visual action recognition and event classification, and achieved strong performance by modeling modality reliability, selecting reliable modalities, performing asymmetric adaptation, or correcting supervision \citep{guo2025smoothing,yang2024read,chen2025tsa,he2026dasp}.
However, existing MM-CTTA methods still mainly rely on backpropagation, typically adapting model parameters through auxiliary objectives such as entropy minimization.
This creates a clear mismatch with practical edge-side deployment: edge devices usually provide limited computation and memory for online adaptation, and deployed models may run with inference engines or quantized components that do not support backward computation.
Therefore, this work studies edge-side MM-CTTA under resource-constrained, forward-only, single-sample streaming deployment, aiming to achieve stable and effective multimodal adaptation without backpropagation.

% 中文：现有不依赖反向传播的适配方法大多基于单模态模型设计。
% 例如，BN-Adapt 主要通过更新 BatchNorm 统计量进行适配，FOA 则依赖 ViT 架构中的 prompt 更新。
% 这些方法的适配接口通常绑定在特定组件上，难以直接迁移到 MM-CTTA 任务中。
% 相比之下，零阶优化（ZOO）提供了一种更通用的前向式学习机制：
% 它通过前向扰动和有限差分估计更新方向，无需反向传播，也不依赖特定模型结构。
% 然而，直接将 ZOO 应用于端侧单样本多模态 CTTA 仍然效果不佳。
% 我们的诊断分析表明，问题主要来自两个方面。
% 第一，高维多模态参数空间会显著降低 ZOO 的搜索效率。
% 多模态模型通常包含模态特定分支和融合模块，相比单模态模型具有更大且更耦合的参数空间。
% 若直接在完整参数空间中进行前向扰动，不仅会降低 ZOO 的收敛速度，
% 还会使有限的前向查询预算分散到大量与当前域偏移弱相关的参数上，
% 导致方向估计的方差增大、信噪比降低，模型难以快速定位真正有效的更新维度。
% 第二，不稳定的局部反馈与熵驱动目标会共同导致更新方向漂移。
% ZOO 通过扰动前后的目标差值近似更新方向，其在线优化过程高度依赖当前测试样本提供的局部反馈。
% 在单样本流式多模态适配中，模态失衡与融合偏置容易使这种反馈变得不可靠。
% 不仅如此，在持续域偏移场景下，基于熵的损失会随样本状态和偏移强度发生明显波动，
% 进一步放大方向估计中的噪声。
% 因此，ZOO 难以获得稳定的学习信号，并容易被瞬时预测偏置牵引，导致适配方向持续漂移。
Existing backpropagation-free adaptation methods are mostly designed for single-modal models.
For example, BN-Adapt mainly adapts models by updating BatchNorm statistics, while FOA relies on prompt updates in ViT architectures \citep{nado2020predictiontime,niu2024forwardonly}.
Their adaptation interfaces are usually tied to specific model components, making them difficult to directly transfer to MM-CTTA.
In contrast, zeroth-order optimization (ZOO) provides a more general forward-only learning mechanism.
It estimates update directions through forward perturbations and finite differences, requires no backpropagation, and does not rely on specific model structures.
However, directly applying ZOO to edge-side single-sample multimodal CTTA remains ineffective in practice.
Our diagnostic analysis identifies two main obstacles.
\textbf{First, high-dimensional multimodal parameter spaces significantly reduce the search efficiency of ZOO.}
Multimodal models usually contain modality-specific branches and fusion modules, resulting in a larger and more coupled parameter space than single-modal models.
If ZOO directly perturbs the full parameter space, it not only slows convergence, but also disperses the limited forward-query budget over many parameters that are only weakly related to the current domain shift.
This increases the variance of direction estimation, reduces the signal-to-noise ratio, and makes it difficult for the model to quickly locate truly effective update dimensions.
\textbf{Second, unstable local feedback and entropy-driven objectives jointly cause update-direction drift.}
ZOO approximates update directions from perturbed objective differences, so its online optimization process heavily depends on the local feedback provided by the current test sample.
In single-sample streaming multimodal adaptation, modality imbalance and fusion bias can easily make such feedback unreliable.
Moreover, under continual domain shifts, entropy-based losses can fluctuate noticeably with sample states and shift severity, further amplifying the noise in direction estimation.
Therefore, ZOO struggles to obtain stable learning signals and can be pulled by instantaneous prediction bias, causing the adaptation direction to drift over time.

% 中文：为解决上述挑战，我们提出 Selective Zeroth-Order Adaptation（SZOA）。
% SZOA 通过选择更有效的更新层，并为 ZOO 构造经过可靠性校准的学习信号，
% 使模型能够在端侧多模态 CTTA 场景中实现稳定的前向式适配。
% 整体而言，SZOA 由两个关键组件组成：Shift-Sensitive Layer Selection（SSLS）
% 和 Reliability-Calibrated Learning Signal Construction（RCLS）。
% 具体而言，SZOA 首先利用 SSLS 从少量 source reference samples 和 early target samples 中
% 提取各模态分支的逐层特征，并通过二类聚类得到的 clustering purity 衡量每一层的域可分性。
% 较高的聚类纯度表示该层仍保留更显著的 source-target 域差异，
% 因此更可能对应当前域偏移的关键适配位置。
% 基于该度量，SZOA 选择 shift-sensitive 层作为 ZOO 的更新空间，
% 从而避免在完整多模态参数空间中进行低效搜索。
% 随后，SZOA 通过 RCLS 在局部流式窗口中估计样本可靠性。
% 具体地，RCLS 依据当前样本预测类别在窗口中的出现频率刻画其预测偏置水平：
% 若该类别在窗口中过度频繁出现，则说明当前样本更可能受到预测偏置影响，因而潜在不可靠；
% 反之，则表明其相对更可靠。
% 在此基础上，RCLS 对样本进行可靠性加权，使低偏置且高置信样本提供正向适配信号，
% 而高偏置或低置信样本被显式抑制，从而增强对可靠样本的学习并减少对不可靠样本的依赖。
% 进一步地，考虑到熵驱动信号在持续域偏移下仍可能剧烈波动，
% RCLS 在 shift-sensitive 层上引入源锚定分布对齐约束，为 ZOO 提供更平滑、更稳定的优化信号。
% 由此，SZOA 同时稳定了 ZOO 的更新位置和学习信号，
% 使其能够在资源受限的单样本多模态流式部署场景中实现高效、稳健的前向式在线适配。
% 这项工作的贡献可以总结如下：
To address the above challenges, we propose a novel
\textbf{S}elective \textbf{Z}eroth-\textbf{O}rder \textbf{A}daptation \textbf{(SZOA)} method, which selects effective update layers and constructs reliability-calibrated learning signals for ZOO, enabling stable forward-only adaptation in edge-side multimodal CTTA.
Overall, SZOA consists of two key components: Shift-Sensitive Layer Selection (SSLS) and Reliability-Calibrated Learning Signal Construction (RCLS).
Specifically, SZOA first uses SSLS to extract layer-wise features from a small number of source reference samples and early target samples in each modality branch, and measures the domain separability of each layer by clustering purity obtained from binary clustering.
Higher clustering purity indicates that a layer still preserves more significant source-target domain discrepancy, and is therefore more likely to correspond to a key adaptation location under the current domain shift.
Based on this metric, SZOA selects shift-sensitive layers as the ZOO update space, avoiding inefficient search over the full multimodal parameter space.
Then, SZOA uses RCLS to estimate sample reliability within a local streaming window.
RCLS characterizes the prediction bias of the current sample according to the frequency of its predicted class in the window: if the class appears overly frequently, the current sample is more likely to be affected by prediction bias and is therefore potentially unreliable; otherwise, it is relatively more reliable.
Based on this estimation, RCLS performs reliability-aware weighting, allowing low-bias and high-confidence samples to provide positive adaptation signals while explicitly suppressing high-bias or low-confidence samples, thereby strengthening learning from reliable samples and reducing dependence on unreliable feedback.
Furthermore, since entropy-driven signals can still fluctuate sharply under continual domain shifts, RCLS introduces a source-anchored distribution alignment constraint on the shift-sensitive layers, providing a smoother and more stable optimization signal for ZOO.
In this way, SZOA stabilizes both the update location and the learning signal of ZOO, enabling efficient and robust forward-only online adaptation in resource-constrained single-sample multimodal streaming deployment.
The contributions of this work are summarized as follows:

% 中文：本文贡献如下。
% 第一，我们提出 SZOA，一种面向端侧多模态持续测试时适配的无反向传播框架。
% SZOA 基于零阶优化，仅依赖前向传播即可在单样本流式测试过程中进行在线适配。
% 第二，我们系统分析了直接将 ZOO 用于多模态 CTTA 时的两个核心障碍：
% 局部前向反馈与熵目标波动导致的方向漂移，以及高维多模态参数空间导致的低效搜索。
% 为此，我们提出 SSLS 和 RCLS，分别从搜索空间和学习信号两个层面稳定 ZOO。
% 第三，我们在 Kinetics50-C 和 VGGSound-C 等多模态 corruption benchmark 上进行系统实验，
% 并在量化 backbone 上进一步验证 SZOA 的端侧部署价值。
\begin{itemize}
    \item We propose SZOA, a backpropagation-free framework for edge-side multimodal continual test-time adaptation. Built on zeroth-order optimization, SZOA performs online adaptation using only forward passes under single-sample streaming deployment.
    \item We provide a diagnostic analysis identifying two core obstacles of applying ZOO to multimodal CTTA: direction drift caused by local forward feedback and entropy fluctuations, and inefficient search caused by high-dimensional multimodal parameter spaces. We address them with SSLS and RCLS, which stabilize ZOO from the search-space and learning-signal perspectives, respectively.
    \item We conduct comprehensive experiments with CAV-MAE on Kinetics50-C and VGGSound-C multimodal corruption benchmarks, and further evaluate SZOA on quantized backbones to validate its edge-side deployment value. Results show that SZOA improves adaptation stability and overall performance while preserving backpropagation-free and low-resource adaptation.
\end{itemize}
